\section{Cycle word algebra}

This chapter develops the algebra of the cycle word: the recursive
construction, the prefix equalities, the closedness equation, and the
QB-8 split.

\subsection{The cycle word}

For a starting depth $c$ and a period length $p$, define
\begin{equation*}
  \mathrm{cycleWord}(c,0)=[],
\end{equation*}
\begin{equation*}
  \mathrm{cycleWord}(c,p+1)
    =\mathrm{cycleWord}(c,p)
    ++ \left[\vtwo{5\,\Orbit_7(c+p)+1}\right].
\end{equation*}
The appended entry is the exact step weight at depth $c+p$.

\subsection{Prefix orbit equality}

\begin{lemma}
For every $c,p$,
\begin{equation*}
  \wordOrbit(\mathrm{cycleWord}(c,p),\Orbit_7(c))
    =\Orbit_7(c+p).
\end{equation*}
\end{lemma}

\begin{proof}
Proceed by induction on $p$.  For $p=0$, the word is empty and both
sides equal $\Orbit_7(c)$.  For $p+1$, write
\begin{equation*}
  w=\mathrm{cycleWord}(c,p),
  \qquad
  t=\vtwo{5\,\Orbit_7(c+p)+1}.
\end{equation*}
By the induction hypothesis,
\begin{equation*}
  \wordOrbit(w,\Orbit_7(c))=\Orbit_7(c+p).
\end{equation*}
Appending $t$ gives
\begin{equation*}
  \wordOrbit(w{++}[t],\Orbit_7(c))
    =\frac{5\,\Orbit_7(c+p)+1}{2^t}
    =\Orbit_7(c+p+1),
\end{equation*}
by the definition of the accelerated step.
\end{proof}

The Lean statement
\texttt{cycleWord\_\allowbreak prefix\_\allowbreak orbit\_\allowbreak eq}
records this equality.

\subsection{Exact step weights}

\begin{lemma}
For every $k<p$,
\begin{equation*}
  \vtwo{5\,\wordOrbit(\mathrm{cycleWord}(c,p){\upharpoonright}k,
  \Orbit_7(c))+1}
  =\mathrm{cycleWord}(c,p)_k.
\end{equation*}
\end{lemma}

\begin{proof}
By the prefix equality, the state after $k$ steps is
$\Orbit_7(c+k)$.  The $k$-th entry of the cycle word was defined to be
the exact valuation at that state.
\end{proof}

This is \texttt{cycleWord\_\allowbreak step\_\allowbreak exact}.  The
companion statement \texttt{cycleWord\_\allowbreak take\_\allowbreak eq}
records the prefix equality for every truncation of the cycle word.

\subsection{Closedness}

Assume $\OrbitCycle(7)$.  There are indices $m<n$ with
\begin{equation*}
  \Orbit_7(m)=\Orbit_7(n).
\end{equation*}
Put $c=m$ and $p=n-m$.  The prefix equality gives
\begin{equation*}
  \wordOrbit(\mathrm{cycleWord}(c,p),\Orbit_7(c))=\Orbit_7(n)
  =\Orbit_7(c),
\end{equation*}
so the cycle word is closed.

\subsection{The cycle equation}

Let $w=\mathrm{cycleWord}(c,p)$, let $m_0=\Orbit_7(c)$, and let
$T=\wordWeight(w)$.  The word-orbit identity gives
\begin{equation*}
  2^T\,m_0=5^p\,m_0+\wordA(w).
\end{equation*}
Closedness supplies the left-hand side through the exact step
division; the identity itself is the compiled statement
\texttt{cycleWord\_\allowbreak cycle\_\allowbreak equation}.

From the equation,
\begin{equation*}
  2^T=5^p+\frac{\wordA(w)}{m_0}>5^p,
\end{equation*}
so
\begin{equation*}
  T>p\log_2 5.
\end{equation*}
This strict inequality is used by the QB-8 split below.

\subsection{The rise/C3 filters}

Define
\begin{equation*}
  \mathrm{rise}=\{t\in w:t<3\},
  \qquad
  \mathrm{c3}=\{t\in w:t\ge3\}.
\end{equation*}
Every rise entry is $1$ or $2$, every C3 entry is at least $3$, and
the two filters partition the word and its total weight:
\begin{equation*}
  |\mathrm{rise}|+|\mathrm{c3}|=p,
  \qquad
  \wordWeight(\mathrm{rise})+\wordWeight(\mathrm{c3})=T.
\end{equation*}

\subsection{Both filters are nonempty}

\begin{proposition}
For a closed cycle word, $\mathrm{rise}$ and $\mathrm{c3}$ are both
nonempty.
\end{proposition}

\begin{proof}
If $\mathrm{c3}$ were empty, then every entry would be at most $2$,
so $T\le2p$.  But the cycle equation gives $T>p\log_2 5$, and
$\log_2 5>2$, contradiction.

If $\mathrm{rise}$ were empty, then every entry would be at least
$3$.  For every positive state $x$ and every $t\ge3$,
\begin{equation*}
  \frac{5x+1}{2^t}\le\frac{5x+1}{8}<x,
\end{equation*}
because $5x+1<8x$ for $x\ge1$.  Hence every step would strictly
decrease the state, and the word could not be closed.
\end{proof}

The Lean records are \texttt{cycleQb8Input\_\allowbreak
hrise\_\allowbreak pos}, \texttt{cycleQb8Input\_\allowbreak
hc3\_\allowbreak pos}, and
\texttt{cycleQb8Input\_\allowbreak P\_\allowbreak ge\_\allowbreak two}.

\subsection{Rise residues modulo 5}

A rise step of weight $1$ satisfies
\begin{equation*}
  2y=5x+1,
\end{equation*}
so modulo $5$,
\begin{equation*}
  2y\equiv1\pmod5,
  \qquad
  y\equiv3\pmod5.
\end{equation*}
A rise step of weight $2$ satisfies
\begin{equation*}
  4y=5x+1,
\end{equation*}
so
\begin{equation*}
  4y\equiv1\pmod5,
  \qquad
  y\equiv4\pmod5.
\end{equation*}
These are the only residue constraints used at the block level; the
Lean statement is
\texttt{rise\_\allowbreak step\_\allowbreak next\_\allowbreak
mod\_\allowbreak five}.

\subsection{The QB-8 input record}

The Lean structure \texttt{CycleQb8Input m S P w rise c3} records the
cycle word together with its filter split, the exact step equations,
the closedness, and the positivity and residue data.  It is produced
from $\OrbitCycle(7)$ by
\texttt{orbit\_\allowbreak cycle\_\allowbreak imp\_\allowbreak
cycle\_\allowbreak qb8\_\allowbreak input}.  The extraction is
compiled.

\subsection{Statement shapes}

\begin{verbatim}
theorem cycleWord_prefix_orbit_eq (c p : Nat) :
    wordOrbit (cycleWord c p) (Orbit 7 c) = Orbit 7 (c + p)

theorem cycleWord_step_exact (c p k : Nat) (hk : k < p) :
    twoValuation
      (5 * wordOrbit ((cycleWord c p).take k) (Orbit 7 c) + 1)
      = (cycleWord c p).getI k

theorem cycleWord_cycle_equation (c p : Nat)
    (hclosed : wordOrbit (cycleWord c p) (Orbit 7 c) = Orbit 7 c) :
    2 ^ wordWeight (cycleWord c p) * Orbit 7 c
      = 5 ^ p * Orbit 7 c + wordA (cycleWord c p)
\end{verbatim}

\subsection{Status}

\begin{center}
\small
\begin{tabular}{@{}p{0.56\textwidth}p{0.36\textwidth}@{}}
\toprule
Statement & Status \\
\midrule
cycle-word recursion & compiled \\
prefix orbit equality & compiled \\
exact step weights & compiled \\
cycle equation & compiled \\
QB-8 split and nonemptiness & compiled \\
rise residue rules & compiled \\
cycle-word extraction from $\OrbitCycle(7)$ & compiled \\
\bottomrule
\end{tabular}
\end{center}
